% Torsion in Poincaré Gauge Theory
% Source: docs/torsion.md (migrated to LaTeX)
% Uses macros from preamble.tex

\section{Torsion in Poincar\'e Gauge Theory}
\label{sec:torsion}

\subsection{Overview}

TIDAL supports Poincar\'e gauge theory (PGT) Lagrangians with propagating torsion. The torsion tensor $\torsion^a{}_{bc}$ is the field strength of the translational gauge symmetry, introduced via xAct's \texttt{DefCovD} with \texttt{Torsion -> True}. Torsion components are perturbed alongside the metric perturbation using \texttt{DefTensorPerturbation}, enabling the full graviton-torsion mixing problem.

\subsection{PGT Framework}
\label{sec:torsion-pgt}

In Poincar\'e gauge theory, gravity is a gauge theory of the Poincar\'e group with two independent gauge fields:

\begin{itemize}
  \item \textbf{Tetrad} (vierbein) $e^a{}_\mu$: gauges translations $\to$ encodes the metric via $\gmunu = \eta_{ab}\, e^a{}_\mu\, e^b{}_\nu$
  \item \textbf{Spin connection} $\omega^{ab}{}_\mu$: gauges Lorentz rotations
\end{itemize}

The corresponding field strengths are:

\begin{itemize}
  \item \textbf{Torsion}: $\torsion^a{}_{\mu\nu} = \partial_\mu e^a{}_\nu - \partial_\nu e^a{}_\mu + \omega^a{}_{b\mu}\, e^b{}_\nu - \omega^a{}_{b\nu}\, e^b{}_\mu$
  \item \textbf{Curvature}: $R^{ab}{}_{\mu\nu}$ (standard Riemann-Cartan curvature)
\end{itemize}

The Riemann-Cartan Ricci scalar $\ricci$ includes torsion contributions and is related to the Levi-Civita Ricci scalar $R$ by the identity:
\begin{equation}
  \ricci = R - \tfrac{1}{4}\, \torsion_{abc}\,\torsion^{abc}
         - \tfrac{1}{2}\, \torsion_{abc}\,\torsion^{bac}
         + \torsion^a{}_{ab}\,\torsion^{cb}{}_{c}
         + \text{total derivative}
  \label{eq:riemann-cartan-identity}
\end{equation}
(Shapiro~\cite{shapiro2002torsion}, eq.~2.17; Hehl et al., 1976)

\subsection{TOML Configuration}
\label{sec:torsion-toml}

\subsubsection{Enabling Torsion}
\label{sec:enabling-torsion}

Add a \texttt{[torsion]} section to the TOML config. This is an optional, standalone section---existing theories without torsion are unaffected.

\begin{lstlisting}[style=toml]
[spacetime]
dimension = 4
metric = "minkowski"

[torsion]
perturbation_name = "t"     # name for the linearized torsion field in equations/JSON
# background = "Tbar"       # future: non-zero background torsion
# irreducible = "axial"     # future: restrict to specific torsion sector
\end{lstlisting}

The \texttt{[torsion]} section:
\begin{enumerate}
  \item Extends the connection: defines \texttt{CDT} with \texttt{DefCovD[..., Torsion -> True]}
  \item Auto-creates \texttt{TorsionCDT[a, -b, -c]} (antisymmetric in \texttt{-b}, \texttt{-c}) for use in the Lagrangian
  \item Registers a perturbation field named by \texttt{perturbation\_name} (e.g., \texttt{t\_0}, \texttt{t\_1}, \ldots, \texttt{t\_8} in the JSON)
  \item Handles background torsion zeroing (flat Minkowski) and xPert label management
\end{enumerate}

\paragraph{Architecture.}
\label{par:torsion-architecture}
The \texttt{[torsion]} section is separate from \texttt{[spacetime]} because in PGT, the metric (from tetrad) and torsion (from connection) are independent gauge fields. The \texttt{[spacetime]} section defines the metric and coordinates; \texttt{[torsion]} extends the connection.

\paragraph{Reserved names.}
\label{par:reserved-names}
The \texttt{perturbation\_name} must not collide with field names, constant names, or built-in operators (\texttt{CD}, \texttt{CDT}, \texttt{eta}, \texttt{Ricci}, \texttt{Torsion}, etc.). The validation system checks this automatically.

\subsubsection{Writing the Lagrangian}
\label{sec:torsion-lagrangian}

The Lagrangian must be written in \textbf{decomposed form} using \texttt{RicciScalarCD[]} (Levi-Civita, NOT \texttt{RicciScalarCDT[]}) plus explicit \texttt{TorsionCDT[...]} terms:

\begin{lstlisting}[style=toml]
[lagrangian]
expression = "(1/kappa^2) RicciScalarCD[] + alpha1 TorsionCDT[-a,-b,-c] eta[a,d] eta[b,e] eta[c,f] TorsionCDT[-d,-e,-f] + alpha2 TorsionCDT[-a,-b,-c] eta[b,d] eta[a,e] eta[c,f] TorsionCDT[-d,-e,-f] + alpha3 TorsionCDT[b,-b,-a] eta[a,d] TorsionCDT[c,-c,-d]"
\end{lstlisting}

\paragraph{Why decomposed form?}
\label{par:decomposed-form}
xPert's built-in perturbation formulas assume the Levi-Civita connection. Using \texttt{RicciScalarCDT[]} directly would require decomposing it via \texttt{ChangeCurvature}, but \texttt{ChangeTorsion} converts torsion $\to$ contortion (a derived quantity that \texttt{DefTensorPerturbation} cannot perturb). The decomposed form avoids this issue entirely.

The three $\torsion^2$ invariants correspond to the three irreducible pieces of torsion under $\mathrm{SO}(1,3)$:
\begin{itemize}
  \item $\torsion_{abc}\,\torsion^{abc}$: couples to all three irreducible sectors
  \item $\torsion_{abc}\,\torsion^{bac}$: the ``mixed'' contraction
  \item $\torsion^a{}_{ab}\,\torsion^{cb}{}_{c}$: the trace sector
\end{itemize}

\subsubsection{Linearization}
\label{sec:torsion-linearization}

\begin{lstlisting}[style=toml]
[linearization]
perturbation_field = "h"   # Metric perturbation (standard xPert)
\end{lstlisting}

Torsion perturbation is \textbf{auto-registered} when a \texttt{[torsion]} section is present. No \texttt{[[linearization.matter\_perturbations]]} entry is needed---the pipeline automatically:
\begin{enumerate}
  \item Defines the perturbation field (e.g., $t_{a,-b,-c}$) from \texttt{perturbation\_name}
  \item Calls \texttt{DefTensorPerturbation} to connect it to \texttt{TorsionCDT}
  \item Handles the xPert label/truncation (see below)
  \item Sets background torsion and contortion to zero (flat Minkowski)
  \item Calls \texttt{VarD} w.r.t.\ the torsion perturbation field for the torsion EOM
\end{enumerate}

\subsubsection{Perturbation Truncation Logic}
\label{sec:perturbation-truncation}

xPert expands each field as: $\Phi = \bg{\Phi} + \varepsilon\,\pert{\Phi}^{(1)} + \varepsilon^2\,\pert{\Phi}^{(2)} + \cdots$
\begin{itemize}
  \item \texttt{LI[1]} = $\pert{\Phi}^{(1)}$---the linear perturbation (dynamical field)
  \item \texttt{LI[2]} = $\pert{\Phi}^{(2)}$---2nd-order correction to the field itself (NOT a separate DOF)
\end{itemize}

The second-order Lagrangian $\lag^{(2)} = \tfrac{1}{2}\,\delta^2 S$ contains three types of terms:
\begin{enumerate}
  \item $(\pert{\Phi}^{(1)})^2$---quadratic products of 1st-order perturbations $\to$ \textbf{KEPT} (give the linear EOM)
  \item $\bg{\Phi} \cdot \pert{\Phi}^{(2)}$---background $\times$ 2nd-order correction $\to$ \textbf{ZERO} (background $\bg{\Phi} = 0$)
  \item $\pert{\Phi}^{(2)}$ alone---appears in $\lag^{(1)}$ only, not $\lag^{(2)}$
\end{enumerate}

Therefore $\texttt{LI[2]} \to 0$ is correct: the physical content is entirely in the \texttt{LI[1]} terms.
This applies identically to metric ($h$), torsion ($t$), and matter ($a$) perturbations.

\subsubsection{Background Contortion Zeroing}
\label{sec:contortion-zeroing}

The contortion $\contorsion^a{}_{bc} = \tfrac{1}{2}(\torsion^a{}_{bc} + \torsion_b{}^a{}_c - \torsion_{bc}{}^a)$ is algebraically determined by torsion. For zero background torsion ($\bg{\torsion} = 0$), the background contortion $\bg{\contorsion} = 0$. In xPert's expansion:
\begin{itemize}
  \item \texttt{Perturbation[K, 0]} = $\bg{\contorsion}$ = 0 $\to$ zeroed via \texttt{ChristoffelCDCDT[\_\_] :> 0}
  \item \texttt{Perturbation[K, 1]} = $\pert{\contorsion}(t)$ $\to$ \textbf{preserved} (encodes the torsion wave)
  \item \texttt{Perturbation[K, 2]} $\to$ already dropped by $\texttt{LI[2]} \to 0$
\end{itemize}

This is exactly analogous to zeroing \texttt{RicciScalarCD[]} (flat background curvature) while keeping \texttt{Perturbation[RicciScalarCD[]]} (the graviton wave).

\subsection{Torsion Irreducible Decomposition}
\label{sec:torsion-irreducible}

Torsion $\torsion_{abc}$ (24 components in 4D) decomposes into three irreducible pieces under $\mathrm{SO}(1,3)$:

\begin{table}[htbp]
\centering
\caption{Irreducible decomposition of the torsion tensor in 4D.}
\label{tab:torsion-irreducible}
\begin{tabular}{@{}llcl@{}}
\toprule
Piece & Name & Components & Description \\
\midrule
${}^{(1)}\torsion_{abc}$ & Tensor (tentor) & 16 & Traceless, tensor part \\
${}^{(2)}\torsion_a = \torsion^b{}_{ba}$ & Trace (trator) & 4 & Vector trace \\
${}^{(3)}\torsion_a = \varepsilon_{abcd}\,\torsion^{bcd}$ & Axial (axitor) & 4 & Totally antisymmetric $\to$ pseudoscalar \\
\bottomrule
\end{tabular}
\end{table}

The PGT Lagrangian parameters $\alpha_1$, $\alpha_2$, $\alpha_3$ control the mass and kinetic terms for each sector independently.

\subsubsection{Propagating Modes}
\label{sec:propagating-modes}

\begin{table}[htbp]
\centering
\caption{Propagating torsion degrees of freedom for different PGT Lagrangians.}
\label{tab:propagating-modes}
\begin{tabular}{@{}ll@{}}
\toprule
PGT Lagrangian & Propagating torsion DOF \\
\midrule
$R$ only (Einstein-Cartan) & 0---torsion is non-propagating \\
$R + \alpha_1 \torsion^2$ & Massive tensor + other modes \\
$R + \alpha_3 \torsion^2_\text{trace}$ & 1 pseudoscalar (axial) \\
$R + \alpha_1 \torsion^2 + \alpha_2 \torsion^2 + \alpha_3 \torsion^2$ & Up to 5 modes \\
\bottomrule
\end{tabular}
\end{table}

\textbf{Key result}: Pure Einstein-Cartan ($\lag = R$ with torsion, no $\torsion^2$ terms) gives NO torsion-Gertsenshtein effect---torsion is non-propagating and vanishes for spinless EM (theorem, not approximation). The $\torsion^2$ kinetic terms are essential for propagating torsion.

\subsection{Implementation Architecture}
\label{sec:torsion-implementation}

\subsubsection{What Happens During Derivation}
\label{sec:torsion-derivation-steps}

\begin{enumerate}
  \item \texttt{[torsion]} section present $\to$ \texttt{DefCovD[CDT, Torsion -> True, FromMetric -> metric]}
  \item Perturbation field defined: \texttt{DefTensor[\{prefix\}T[a,-b,-c], M, Antisymmetric[\{-b,-c\}]]}
  \item \texttt{DefTensorPerturbation} connects \texttt{\{prefix\}tPert[LI[order],...]} to \texttt{TorsionCDT[...]}
  \item Lagrangian parsed with \texttt{RicciScalarCD[]} + \texttt{TorsionCDT[...]} terms
  \item \texttt{SetupMetricPerturbation} for metric $h$
  \item xPert: $\lag^{(2)}$ = second-order perturbation in both $h$ and torsion
  \item Truncation: \texttt{tPert[LI[2],...] $\to$ 0}, \texttt{tPert[LI[1],...] $\to$ T[...]}, \texttt{TorsionCDT[...] $\to$ 0}
  \item \texttt{VarD}: derive EOM for both $h$ and $\torsion$ components
  \item Component decomposition $\to$ plane-wave reduction $\to$ constraint elimination
\end{enumerate}

\subsubsection{Why ChangeCurvature/ChangeTorsion Were NOT Used}
\label{sec:why-not-changecurvature}

Investigation revealed:
\begin{itemize}
  \item \texttt{ChangeCurvature[L, CDT, CD]} correctly decomposes $\ricci \to R$ + contortion terms
  \item But the contortion terms involve \texttt{Christoffel[CD,CDT]}---a derived xAct object
  \item \texttt{ChangeTorsion[L, CDT, CD]} converts \texttt{TorsionCDT} $\to$ \texttt{Christoffel[CD,CDT]} (the WRONG direction)
  \item \texttt{DefTensorPerturbation} requires fundamental tensors, not derived Christoffel differences
  \item xAct does not provide the inverse conversion (contortion $\to$ torsion) natively
\end{itemize}

The decomposed-form Lagrangian avoids this issue entirely.

\subsection{References}
\label{sec:torsion-references}

\begin{description}
  \item[Blagojevi\'c \& Hehl (2013)] \emph{Gauge Theories of Gravitation}. Imperial College Press.
  \item[Hehl et al.\ (1976)] ``General Relativity with Spin and Torsion.'' \emph{Rev.\ Mod.\ Phys.}\ 48:393. DOI: \href{https://doi.org/10.1103/RevModPhys.48.393}{10.1103/RevModPhys.48.393}
  \item[Shapiro (2002)~\cite{shapiro2002torsion}] ``Physical Aspects of the Space-Time Torsion.'' \emph{Phys.\ Rept.}\ 357:113. arXiv: \href{https://arxiv.org/abs/hep-th/0103093}{hep-th/0103093}
  \item[Hayashi \& Shirafuji (1979)] ``New General Relativity.'' \emph{Phys.\ Rev.\ D}\ 19:3524.
  \item[Nikiforova et al.\ (2009)] ``Stability of the Massive Torsion Modes.'' arXiv: \href{https://arxiv.org/abs/0905.4007}{0905.4007}
  \item[Barker (2023)~\cite{barker2024poincare}] ``HiGGS: Hamiltonian in Gauge Gravity Solver.'' \emph{Eur.\ Phys.\ J.\ C}\ 83:228. arXiv: \href{https://arxiv.org/abs/2206.00658}{2206.00658}
  \item[xAct] \href{http://www.xact.es/Documentation/HTML/HTMLLinks/xTensor/DefCovD.nb.html}{DefCovD with Torsion}~\cite{martingarcia2007xact}.
\end{description}
